\documentclass[12pt,a4paper]{article}
% Estilo compartido para informes. Compilar con LuaLaTeX desde tp1/.
\usepackage[spanish,es-nodecimaldot]{babel}
\usepackage{fontspec}
\setmainfont{Montserrat-Medium.ttf}[
  Path=../assets/fonts/,
  BoldFont=Montserrat-Bold.ttf,
  ItalicFont=Montserrat-MediumItalic.ttf]
\usepackage[a4paper,margin=2.5cm,headheight=16pt]{geometry}
\usepackage{xcolor,graphicx,booktabs,tabularx}
\usepackage{float}
\usepackage{titlesec,fancyhdr,enumitem,caption}
\usepackage{ragged2e}
\usepackage{hyperref}
\definecolor{verde}{HTML}{1D4036}
\definecolor{bronce}{HTML}{D38129}
\definecolor{rosa}{HTML}{D88E77}
\definecolor{calido}{HTML}{EFE4D8}
\definecolor{grafito}{HTML}{1D1D1D}
\hypersetup{colorlinks=true,linkcolor=verde,urlcolor=verde,
  pdftitle={TP1 - Instalación de máquinas virtuales},pdfsubject={Cyberseguridad}}
\AtBeginDocument{\color{grafito}\RaggedRight\fontsize{12}{14.4}\selectfont}
\setlength{\parindent}{0pt}
\setlength{\parskip}{7pt}
\setlist[itemize]{leftmargin=*,itemsep=7pt,topsep=4pt}
\setlist[enumerate]{leftmargin=*,itemsep=6pt}
\titleformat{\section}{\color{verde}\bfseries\fontsize{18}{21.6}\selectfont}{\thesection}{0.65em}{}
\titleformat{\subsection}{\color{verde}\bfseries\large}{\thesubsection}{0.65em}{}
\titleformat{\subsubsection}{\color{verde}\bfseries}{\thesubsubsection}{0.65em}{}
\titlespacing*{\section}{0pt}{16pt}{8pt}
\titlespacing*{\subsection}{0pt}{12pt}{5pt}
\titlespacing*{\subsubsection}{0pt}{10pt}{4pt}
\setcounter{tocdepth}{2}
\pagestyle{fancy}
\fancyhf{}
\fancyhead[L]{\small\color{verde}FCEFyN | UNC}
\fancyhead[R]{\small\color{verde}\materia\ · TP \numeroTP}
\fancyfoot[L]{\scriptsize Informe de laboratorio}
\fancyfoot[R]{\small\thepage}
\renewcommand{\headrulewidth}{0.4pt}
\captionsetup{font=small,labelfont=bf,justification=raggedright,singlelinecheck=false,hypcap=false}
\newcommand{\pendiente}[1]{\par{\small\color{verde}\textbf{Por completar.} #1\par}}
% Si falta una captura se muestra un espacio reservado; no impide compilar.
\newcommand{\captura}[2]{%
  \par\smallskip\begingroup\centering
  \IfFileExists{figuras/#1}{%
    \includegraphics[width=\linewidth,height=5cm,keepaspectratio]{figuras/#1}%
  }{%
    \setlength{\fboxsep}{10pt}%
    \fcolorbox{bronce}{calido}{\parbox[c][1.25cm][c]{\dimexpr\linewidth-2\fboxsep-2\fboxrule\relax}{%
      \centering\small Captura pendiente\par\texttt{\detokenize{#1}}}}%
  }%
  \captionof{figure}{#2}\par\endgroup\smallskip
}

\usepackage{xurl}

\newcommand{\materia}{Cyberseguridad}
\newcommand{\numeroTP}{15}
\newcommand{\tituloTP}{Análisis de una infraestructura real}
\newcommand{\estudiante}{Brigido Tomas Andres}
\newcommand{\legajo}{38481096}
\newcommand{\carrera}{ICOMP}
\newcommand{\docente}{Javier A. Jorge}


% Metadatos específicos del TP15 y capturas en img/.
\hypersetup{pdftitle={TP15 - Análisis de una infraestructura real},linkcolor=black,urlcolor=black,citecolor=black}
\renewcommand{\captura}[2]{%
  \begin{figure}[H]
    \centering
    \IfFileExists{img/#1}{%
      \includegraphics[width=0.85\linewidth,height=5cm,keepaspectratio]{img/#1}%
    }{%
      \setlength{\fboxsep}{10pt}%
      \fcolorbox{bronce}{calido}{\parbox[c][1.1cm][c]{0.85\linewidth}{%
        \centering\small Captura pendiente\par\texttt{\detokenize{#1}}}}%
    }%
    \caption{#2}
  \end{figure}
}

\usepackage{caption}
\captionsetup{justification=centering}
\begin{document}
\hypersetup{pageanchor=false}
\begin{titlepage}
  \includegraphics[width=8cm]{../assets/logo-fcefyn-2026.png}
  \vspace{1cm}

  {\color{verde}\small UNIVERSIDAD NACIONAL DE CÓRDOBA\par}
  \vspace{2.1cm}
  {\color{bronce}\bfseries\large TRABAJO PRÁCTICO \numeroTP\par}
  \vspace{0.45cm}
  {\color{verde}\bfseries\fontsize{32}{38.4}\selectfont Análisis de una\\infraestructura real\par}
  \vspace{0.3cm}
{\small Seguridad y propuesta de mejoras para una radio\par}
\vspace{0.35cm}
  {\large \materia\par}
  \vspace{0.8cm}
  {\color{bronce}\rule{2.7cm}{2pt}\par}
  \vfill
  \renewcommand{\arraystretch}{1.45}
  \begin{tabular}{@{}p{3.2cm}p{10cm}@{}}
    Estudiante & \estudiante \\
    Legajo & \legajo \\
    Carrera & \carrera \\
    Docente & \docente \\
  \end{tabular}
  \vfill
  {\small Córdoba, Argentina\par}
\end{titlepage}
\hypersetup{pageanchor=true}

\tableofcontents
\vspace{1cm}
\newpage

\section{Introducción y alcance}
El trabajo propone analizar la infraestructura de una radio cooperativa, identificar sus componentes críticos y plantear mejoras viables. Se estudiarán la red, los equipos y las aplicaciones a partir de la documentación suministrada.

\section{Contexto de la organización}
La documentación describe un servidor HPE ProLiant MicroServer Gen10 con un SSD de 256 GB para OpenMediaVault 6.0 y un disco de 1 TB para datos. Se utiliza para compartir archivos y facilitar el trabajo remoto mediante Samba, SFTP, FileBrowser, Nextcloud y Hamachi, con Portainer para administrar contenedores.

La radio combina puestos internos, equipos domiciliarios y dispositivos de asistentes a talleres. El documento indica que la separación de las redes wifi, el control del ancho de banda y los respaldos automáticos todavía no estaban implementados. El streaming se realizaba desde una PC dedicada.

Estos datos corresponden al material de referencia de octubre de 2023; no constituyen una verificación del estado actual. Aunque el croquis muestra dos redes wifi, el texto aclara que entonces existía una sola.

Los activos más críticos son:
\begin{itemize}
    \item \textbf{Servidor y archivos compartidos:} su falla o compromiso podría provocar pérdida, filtración o indisponibilidad de los contenidos de trabajo.
    \item \textbf{PC de streaming y puestos de operación:} su interrupción afectaría la transmisión y la producción de la radio.
    \item \textbf{Routers y conexión a Internet:} su falla impediría el acceso remoto y el streaming; una intrusión podría comprometer las comunicaciones.
    \item \textbf{Credenciales de acceso y administración:} su robo permitiría acceder sin autorización a archivos, servicios y configuraciones.
\end{itemize}

\newpage
\subsection*{Elementos de protección incorporados}
\begin{itemize}
    \item \textbf{VPN Hamachi:} permite el acceso remoto mediante una red privada virtual.
    \item \textbf{SFTP:} cifra las transferencias de archivos.
    \item \textbf{Nextcloud:} dispone de cuentas individuales para los usuarios.
    \item \textbf{FileBrowser:} requiere usuario y contraseña, aunque la cuenta compartida limita el control y la identificación individual.
\end{itemize}

\subsection*{Elementos de protección que podrían incorporarse y viabilidad}
\begin{itemize}
    \item \textbf{Respaldos automáticos con copia externa:} viables si se dispone de almacenamiento adicional y se programan y verifican las copias.
    \item \textbf{Cuentas individuales y permisos mínimos:} viables mediante la configuración de usuarios y accesos en los servicios utilizados.
    \item \textbf{Separación de la red de invitados:} depende de que el equipamiento permita aislarla de la red interna; de lo contrario, requiere adquirir equipos compatibles.
    \item \textbf{Prioridad del streaming:} viable si el router permite gestionar el ancho de banda; si no, requiere equipamiento compatible.
\end{itemize}

\clearpage
\section{Análisis de seguridad}
\subsection{Infraestructura de red y arquitectura}
La red conecta el servidor, los puestos de trabajo y los dispositivos wifi. 

Según la documentación, \textbf{personal} e \textbf{invitados} comparten una misma red inalámbrica y no existe control del ancho de banda. Esto puede facilitar accesos indebidos y afectar el streaming.La conexión a Internet y los equipos de red son críticos.

\par\textbf{¿Que herramientas de autorización y autenticación se podrían implementar?}
\begin{itemize}
    \item \textbf{Credenciales individuales para la VPN:} viable si se dispone de un sistema de gestión de identidades.
    \item \textbf{Autenticación multifactor:} viable si el equipamiento lo permite y se puede integrar con los servicios existentes.
    \item \textbf{Restricción de administración:} viable mediante la configuración de permisos y roles en los dispositivos y servicios.
    \item \textbf{Los invitados deberían acceder únicamente a Internet y a los recursos del taller que necesiten:} viable si se implementan controles de acceso adecuados.
\end{itemize}
\par\textbf{¿Que otras herramientas podrían agregarse para asegurar el sistema?}
\begin{itemize}
    \item \textbf{Separación de redes:} aislar la red interna de la de invitados para proteger los recursos de trabajo.
    \item \textbf{Firewall:} aplicar reglas que permitan únicamente las conexiones necesarias y restrinjan accesos no autorizados.
    \item \textbf{Prioridad del streaming:} gestionar el ancho de banda para evitar que otras actividades afecten la transmisión.
    \item \textbf{Monitoreo especializado:} sería conveniente contratar un equipo o servicio externo de ciberseguridad para analizar el tráfico de red, revisar registros de acceso y eventos, detectar actividad sospechosa y coordinar la respuesta ante incidentes.
\end{itemize}
Las medidas técnicas dependen de las funciones del equipamiento disponible. La contratación del servicio de ciberseguridad debe evaluarse según el presupuesto y las necesidades de la radio.
\par\textbf{¿Que bugs o errores puede tener cada uno de estos dispositivos o aplicaciones?}
\begin{itemize}
    \item \textbf{Firmware desactualizado:} podría mantener fallos de seguridad sin corregir.
    \item \textbf{Contraseñas predeterminadas:} podrían facilitar el acceso no autorizado a la administración de los equipos.
    \item \textbf{Puertos abiertos innecesariamente:} podrían exponer servicios que no necesitan estar accesibles.
    \item \textbf{Reglas de acceso mal configuradas:} podrían permitir conexiones a recursos que deberían estar restringidos.
\end{itemize}

\clearpage
\subsection{Nodos y sistemas operativos}
Los nodos incluyen el servidor, los puestos de operación, edición y oficina, la PC de streaming y los equipos personales. El servidor almacena los archivos compartidos y la PC de streaming sostiene la transmisión. Se documenta OpenMediaVault basado en Debian en el servidor. 
\par\textbf{¿Qué herramientas de autorización y autenticación se podrían implementar?}
\begin{itemize}[itemsep=3pt]
    \item \textbf{Cuentas individuales:} identificar a cada usuario y evitar compartir credenciales.
    \item \textbf{Permisos mínimos:} limitar el acceso a archivos y reservar los privilegios administrativos para quienes los necesiten.
    \item \textbf{Acceso administrativo seguro:} usar claves SSH y un segundo factor donde sea compatible, además del bloqueo de pantalla en los puestos locales.
\end{itemize}

\par\textbf{¿Qué otras herramientas podrían agregarse para asegurar el sistema?}
\begin{itemize}[itemsep=3pt]
    \item \textbf{Actualizaciones y protección local:} aplicar parches, activar el firewall del equipo y utilizar protección antimalware compatible con cada sistema.
    \item \textbf{Respaldos y recuperación:} automatizar copias externas de los archivos críticos y comprobar que puedan restaurarse.
    \item \textbf{Protección física y cifrado:} restringir el acceso al servidor y cifrar los equipos portátiles que almacenen información sensible.
\end{itemize}

\par\textbf{¿Qué bugs o errores podrían tener estos equipos o sistemas?}
\begin{itemize}[itemsep=3pt]
    \item \textbf{Sistemas desactualizados:} podrían conservar vulnerabilidades sin corregir o carecer de soporte.
    \item \textbf{Permisos excesivos:} podrían facilitar modificaciones indebidas y aumentar el impacto de una infección.
    \item \textbf{Servicios innecesarios:} podrían ampliar las posibilidades de acceso no autorizado.
    \item \textbf{Falta de respaldos automáticos:} está documentada y aumenta el riesgo de pérdida de información ante fallos o incidentes.
\end{itemize}
Los primeros tres puntos requieren verificación en los equipos; no se han confirmado vulnerabilidades específicas.

\clearpage
\subsection{Aplicaciones y servicios}
\pendiente{Revisar OpenMediaVault, Samba, SFTP, Portainer, FileBrowser, Nextcloud y Hamachi. Considerar las cuentas compartidas, los privilegios administrativos, el cifrado y la exposición de servicios.}
\pendiente{Identificar componentes críticos y protecciones ya incorporadas según la documentación.}

\par\textbf{¿Que herramientas de autorización y autenticación se podrían implementar?}
\pendiente{Respuesta breve aplicada a esta área.}
\par\textbf{¿Que otras herramientas podrían agregarse para asegurar el sistema?}
\pendiente{Indicar medidas y justificar su viabilidad técnica, económica y operativa.}
\par\textbf{¿Que bugs o errores puede tener cada uno de estos dispositivos o aplicaciones? Ej. versiones de fw desactualizadas, versiones de aplicaciones desactualizadas ... etc.}
\pendiente{Distinguir riesgos potenciales de fallos confirmados. Para citar vulnerabilidades concretas, comprobar modelo, versión y fuente.}

\clearpage
\section{Autodiagnóstico de seguridad}
La consigna propone responder la herramienta de INCIBE como integrantes de la organización: \url{https://adl.incibe.es/}.

\pendiente{Registrar las respuestas y los supuestos utilizados. Si falta información, dejar constancia de ello.}

\begin{tabularx}{\linewidth}{@{}lXX@{}}
\toprule
\textbf{Área} & \textbf{Resultado obtenido} & \textbf{Observación principal} \\
\midrule
Personas & Por completar & Por completar \\
Procesos & Por completar & Por completar \\
Tecnología & Por completar & Por completar \\
\bottomrule
\end{tabularx}

Los archivos Personas\_alto, Procesos\_alto y Tecnologia\_alto se utilizarán como material de apoyo. Su nivel de riesgo no se atribuirá a un nuevo diagnóstico sin haber realizado y documentado el cuestionario.

\section{Propuesta de mejoras y prioridades}
\pendiente{Seleccionar las mejoras más relevantes del análisis y ordenar su implementación. Considerar el presupuesto, el conocimiento técnico disponible y la continuidad de la radio.}

\begin{tabularx}{\linewidth}{@{}p{2cm}XXX@{}}
\toprule
\textbf{Prioridad} & \textbf{Mejora} & \textbf{Riesgo que reduce} & \textbf{Viabilidad} \\
\midrule
Por definir & Por completar & Por completar & Por completar \\
Por definir & Por completar & Por completar & Por completar \\
Por definir & Por completar & Por completar & Por completar \\
\bottomrule
\end{tabularx}

\clearpage
\section{Conclusión}
\pendiente{Resumir los principales riesgos y las mejoras prioritarias, sin extenderse.}

\clearpage
\section{Referencias}
\begin{itemize}
\item Consigna del TP15: análisis de una infraestructura real.
\item Cooperativa de trabajo. Propuesta y descripción del servidor HPE ProLiant Gen10, octubre de 2023.
\item Croquis de infraestructura suministrado por la cátedra.
\item INCIBE. Materiales de autodiagnóstico: Personas, Procesos y Tecnología, nivel alto.
\item INCIBE. Herramienta de autodiagnóstico: \url{https://adl.incibe.es/}.
\end{itemize}
\end{document}
